% Source: https://artofproblemsolving.com/community/c3h1473463_dont_understand
% Stable identity: aops-topic-1473463-q01
\textbf{Counting the terms.} After combining like terms, each term is a monomial $a^pb^qc^rd^s$ with $p+q+r+s\le N$, and every such monomial does occur (with a positive coefficient). We need $p,q,r,s\ge1$. Put $t=N-(p+q+r+s)\ge0$ (the power of $1$). Then the terms correspond to solutions of
\[
p+q+r+s+(t+1)=N+1
\]
in positive integers, and by stars and bars there are $\dbinom{N}{4}$ of them. (The AoPS Wiki sum $\dbinom{N-1}{3}+\dbinom{N-2}{3}+\cdots+\dbinom{3}{3}$ counts the same monomials by their total degree $k=p+q+r+s$: there are $\dbinom{k-1}{3}$ of degree $k$, and the hockey stick identity adds them up to $\dbinom{N}{4}$.)

Now $\dbinom{N}{4}=1001=7\cdot11\cdot13$, and $\dbinom{14}{4}=\dfrac{14\cdot13\cdot12\cdot11}{24}=1001$. Since $\dbinom{N}{4}$ increases with $N$, $N=14$, answer $\boxed{\textbf{(B)}}$.

\textbf{Why the guess is wrong.} With $x=a+b+c+d$, the binomial theorem gives $(x+1)^N=\sum_j\dbinom{N}{j}x^j$. The numbers $\dbinom{N}{j}$ are \emph{coefficients} of the powers $x^j$, not numbers of terms. Each $x^j=(a+b+c+d)^j$ expands into many different monomials (for example $x^2$ already has $10$), and many of them, such as $a^j$, do not contain all four variables. So the guessed sum neither counts monomials nor checks which variables appear.
